\documentclass[12pt]{article}

\usepackage[utf8]{inputenc} % for accented characters and €
\usepackage[T1]{fontenc}    % better font encoding
\usepackage{textcomp}       % for \texteuro
\usepackage{hyperref}       % for \url
\usepackage{geometry}       % nicer margins
\usepackage{parskip}        % adds spacing between paragraphs

\geometry{margin=1in}

\title{Bingoals 2025: My Year in Goals}
\author{Marlon Valcke}
\date{}
\begin{document}

\textbf{New Year's resolutions} --- I've never really participated in
those. They're fleeting, poorly thought out, demotivating, and
unproductive. One big pass/fail moment --- and let's be honest, almost
always a fail already in January --- that doesn't motivate me.

But having goals to work towards, and giving yourself a bit of direction
when you don't know what to do with a free moment, \emph{that} can be
incredibly valuable. So how do I do that in a better way?

That was the problem I went to sleep with somewhere in December 2024,
and I woke up the next morning with the solution: \textbf{BINGOALS}.

This idea provides gradients of success, accomplishments throughout the
year, sustained motivation, variation, and flexibility. And it only
vaguely resembles a game stereotypically played by retirees in a
run-down parish hall where you can barely see the floor through the
nicotine clouds hanging in the air.

\textbf{How does it work?}

Brainstorm New Year's resolutions. Throw them all onto a page --- the
big ones, the small ones, the important ones, and the small fun things.
Is it unlikely? Or something you were almost certainly going to do
anyway? That doesn't matter.

Throw everything onto a page and add or remove until you end up with
\textbf{21--25 items}.

These are your New Year's resolutions for the coming year.

And yes, if \emph{not} achieving a single New Year's resolution is
demotivating\ldots{} how could this unrealistic list possibly be better?

Because we \emph{know} we won't achieve every goal. And that that's not
required to ``win.''

Next to our list of goals, we hang up our \textbf{BINGOAL card} for the
year:\\
a \textbf{5×5 grid}, with empty squares.

But those squares won't stay empty for long.

Every time we achieve a goal from our list, we choose a \textbf{random*}
square and write the title of the goal in it.\\
Will we manage to fill a complete line this year and score a
\textbf{BINGO}?

\textbf{Yes we can!}

With a mix of safe and simple goals, the Bingoal card will already start
to fill in a bit --- but in most cases, the actual win will only happen
later in the year.\\
(On average, your first Bingo happens around the \textbf{15th filled
square}.)

* \emph{Random square}: roll a die to determine the row (reroll on a 6).
Then do the same for the column if there's more than one free square
left. If you land on an already filled square, reroll the last die.

Okay, sounds good on paper --- but how does it work in practice?

Well, luckily I tested that specifically for you (and maybe a little bit
for myself).

In short: \textbf{fantastic}.

It gives structure, a sense of achievement, insight into what you
\emph{think} you find important but in practice actually don't care
about at all --- and most importantly, it's just \emph{fun}.

But maybe the best way to judge that is to look at the goals themselves,
and how I experienced them.

\textbf{Examples of Bingoals 2025}\\
(Some are not listed here because I didn't really do anything with them
and don't have anything interesting to say about them.)

\begin{itemize}
\item
  Run 5 km
\item
  Run 10 km
\item
  Teach my dog agility (5 specific tricks noted)
\item
  Earn €100 with 3D models
\item
  Earn €500 with 3D models and prints
\item
  Write my own board game or RPG
\item
  Don't drink cola for a month
\item
  Do something on a stage
\item
  Read the entire \emph{Wheel of Time} series together with my
  girlfriend
\item
  Read 30 books
\item
  Take a walk in every Belgian province
\item
  Walk part of the Müllerthal Trail
\item
  Alphabet dates
\item
  Climb a mountain
\item
  Go camping
\item
  Moisturize for a month
\item
  Install fly screens
\item
  Watch 10 ``real'' films
\item
  Lose 5 kg
\item
  Run or play 20 sessions with the Resonance Protector's Guild (gaming
  group)
\item
  100 days of Duolingo French
\item
  Do a stereotypical group activity (karting, paintball, etc.)
\end{itemize}

Want to know how each of these were experienced, that takes more space
than I get here but you can check that out at
GM-valcke.github.io/bingoals2025

\textbf{Run 5 km}

A good goal in itself.\\
In the past, I've achieved this multiple times via Start to Run
programs, and this year I started one again using music tracks I made
for that purpose about 10 years ago.

Between March and October, I went running occasionally --- especially in
March, April, and May, often enough to notice progress.

However: a 4-kg dog with short legs and a strong urge to walk \emph{in
front of you} or suddenly stop for a pee without any respect for your
rhythm does not make this easier.

As the running intervals became longer and the weather warmer, it became
clear that this wasn't healthy for my running buddy. The idea of first
running for half an hour and \emph{then} immediately having to walk that
small fluffy creature was absolutely not appealing --- so from summer
onwards, I basically stopped running.

\textbf{Success?} No\\
\textbf{Was this a good goal?} Yes\\
\textbf{Will it be on the list for 2026?} Probably not

\textbf{Run 10 km}

This was a continuation of the previous goal.\\
Achieving this second goal depended entirely on achieving the first ---
which means I should have set one more modest goal and either let
enjoyment carry me further, or move the boundary in a later year.

\textbf{Success?} No\\
\textbf{Was this a good goal?} No\\
\textbf{Will it be on the list for 2026?} No

\textbf{Teach my dog agility (5 specific tricks)}

A very nice goal.\\
However, we didn't actually practice specific agility tricks.

This is a goal where I never really took the step to start. In 2024 I
had a false start, and I think that made me hesitant to begin again
unless I was sure I'd stick with it --- and that confidence just wasn't
there this year.

\textbf{Success?} No\\
\textbf{Was this a good goal?} Yes\\
\textbf{Will it be on the list for 2026?} Probably

\textbf{Earn €100 with 3D models}

This was a goal I personally wrote down as a bit of a fantasy.

And yet, I managed to complete it better than expected.

I didn't really invest in promotion or attention for the designs I made.
I simply uploaded designs that I thought others might find interesting
to Cults3D, put what felt like a fair price on them, and then\ldots{}
mostly nothing happened.

But occasionally, designs \emph{were} bought.

In total, I posted around 20 designs and earned about \textbf{€40}.\\
That's not exactly the goal I set, but I'm still happy with it.

The goal was never meant as a real income model --- it was meant as
motivation to finish designs far enough that they were publicly
shareable, instead of stopping at the first functional prototype.

Whether that worked, you can judge for yourself:\\
\url{https://cults3d.com/en/users/marlonvalck/3d-models}

* \emph{This goal should have been phrased differently.}\\
I can decide to make 20 designs --- I either put in the effort or I
don't.\\
But setting a goal like selling a design 100 times introduces a big luck
factor that's not linked to my effort or core objective.

\textbf{Success?} Yes (technically no, but it felt like a success ---
see next goal)\\
\textbf{Was this a good goal?} Yes, but see note above\\
\textbf{Will it be on the list for 2026?} Likely

\textbf{Earn €500 with 3D models and prints}

Just like the 10 km goal, this was an extension of the previous one ---
and I already knew it was unrealistic. That's why I added making prints
to the goal.

How would that work? No idea. I make prints for friends for free, and
selling 3D flexi dragons at flea markets wasn't exactly the plan.
(Though if you're on my friends list and want one, feel free to message
me.)

Because it was on the list, I did occasionally see opportunities to
spontaneously offer help and solve a small problem for a small fee. If I
count my time, it definitely didn't earn me anything --- but they were
mostly fun little projects.

Did I earn €500? Nope.\\
But I \emph{do} feel like I fulfilled the intention of the goal.

Just like the previous one, this wasn't completed to the letter, but it
\emph{was} completed in spirit --- so I decided to count one of the two
as a success.

\textbf{Success?} No (but also a little bit yes)\\
\textbf{Was this a good goal?} Yes\\
\textbf{Will it be on the list for 2026?} Unlikely

\textbf{Write my own board game or RPG}

Over the past few years, several game ideas have popped into my head.\\
But the step from idea to finished product is long and filled with
boring, tedious steps.

\emph{The Toybox} and \emph{Sleeping Dogs} are two board games that
exist as rough drafts on my computer, but if I'm honest, they'll never
be finished.

Two other projects are further along. One is still a secret, but some of
you have already played it. The other --- \textbf{Spellwrighters} --- is
a fantastic improv-storytelling RPG that has been playtested several
times over the past months.

Once again, coming up with ideas and playtesting was far more fun than
writing the game itself --- but \textbf{Spellwrighters exists}. This
goal was achieved!

Interested in experiencing the ultimate magical power fantasy once
\emph{Spellwrighters} is released?\\
Let me know via \url{https://gm-valcke.github.io/contact/} and I'll
happily put together a mailing list to keep you informed first.

\textbf{Success?} YES\\
\textbf{Was this a good goal?} YES\\
\textbf{Will it be on the list for 2026?} Maybe --- but actually
publishing it, creating a proper website and distribution platform
definitely will be.

\textbf{Don't drink cola for a month}

You'll never see me use alcohol or drugs, but apparently everyone needs
an addiction --- and for me, that's Coca-Cola. If you have to be
addicted to coke, this is the least bad option.

A month without this addiction is a goal I've tried before, and I've
always succeeded. Though not without a certain level of withdrawal
symptoms and grumpiness.

This year I kept postponing it, but in November I finally bit the
bullet.\\
(Or rather: drank a \emph{lot} of apple juice to still get some sugar
water into my system.)

Now that I'm writing this in December, I can say I'm once again enjoying
that cola a bit extra. The goal was never to quit, but to recalibrate
and re-appreciate it --- and in that sense, I definitely consider this a
success.

\textbf{Success?} Yes\\
\textbf{Was this a good goal?} Yes\\
\textbf{Will it be on the list for 2026?} Probably

\textbf{Do something on a stage}

This was a poorly thought-out goal.

It seemed like a fun thing to challenge myself with. I figured an
opportunity would probably come along if I kept my eyes open. Maybe it
would push me to take a few improv theatre classes again? Maybe the
juggling / circus itch would come back? (And to be fair, it did come
back a little this summer at the European Juggling Convention.)

Did I stand in front of large groups? Yes --- but teaching, even to 500
students, wasn't the kind of performance I had in mind.

Did I entertain? Sure --- but despite the saying ``three is a crowd,'' I
don't really count the back of a GMS screen as a ``stage.''

Maybe sharing this post is a way of giving myself a stage through social
media?\\
Who knows --- if this post and the Bingoals concept go viral around New
Year's, I might still tick this one off.

\textbf{Success?} No\\
\textbf{Was this a good goal?} No (but maybe not all Bingoals should be
good ones?)\\
\textbf{Will it be on the list for 2026?} No

\textbf{Read the entire \emph{Wheel of Time} series together with my
girlfriend}

I've already read this series six or seven times. It's the foundation
upon which my love for fantasy and stories is built. Being able to read
it together and share that experience is just incredibly fun.

Convincing your girlfriend to read 1,000 pages of fantasy per month
together?\\
\#relationshipgoals

\textbf{Success?} YES\\
\textbf{Was this a good goal?} YES\\
\textbf{Will it be on the list for 2026?} No --- but a lighter version
of this goal might return

\textbf{Read 30 books}

This is essentially the annual Goodreads challenge. Since the heavy
hitters of \emph{The Wheel of Time} would be included this year, I kept
the number modest.

I think finding good series was especially important to keep the pace
going. A fun goal that will definitely return --- and one I've actually
completed every year since 2016, except for one.

Somewhere in November, after reaching the goal, I started \emph{Arabian
Nights} --- and suddenly my reading came to a halt. Unlike what happens
in the story, I was genuinely not motivated to find out what would
happen next in Scheherazade's tale.

\textbf{Success?} YES\\
\textbf{Was this a good goal?} Yes\\
\textbf{Will it be on the list for 2026?} Yes --- probably with fewer
books, but a bit more challenge in terms of content. I'd like to add a
non-fiction reading challenge, which isn't something I'd normally pick
up.

\textbf{Take a walk in every Belgian province}

Millennials in 2025: vacation destination --- the other side of the
world.

That seems like such a waste when we barely know our own backyard
anymore. Let this be an excuse to explore our small country a bit more.

These were great excuses to head out for a Saturday walk. Was every
destination equally enchanting? No. But each one offered something
unique. And I feel like it helped me better understand the scale of
Belgium.

I'll spare you the full list and stick to the highlights:

\begin{itemize}
\item
  \textbf{Best walk:} Promenade de la Hoëgne (Liège)
\item
  \textbf{Best photo:} squirrel in Molsbroek (Lokeren, East Flanders)
\item
  \textbf{Biggest disappointment:} Villers-la-Ville (Walloon Brabant)
  --- but only because we couldn't walk in the nicest areas due to
  hunters, and right afterwards the dog's collar broke, forcing me to
  carry him --- flailing --- back to the car.
\end{itemize}

\textbf{Success?} Yes\\
\textbf{Was this a good goal?} Yes\\
\textbf{Will it be on the list for 2026?} No

\textbf{Walk part of the Müllerthal Trail}

This ended up on the paper as a ``bonus goal'' for the previous one. I
completely forgot it was a separate goal until far too late in the year.

This one gets a second chance next year.

\textbf{Success?} No\\
\textbf{Was this a good goal?} Yes\\
\textbf{Will it be on the list for 2026?} Yes

\textbf{Alphabet dates}

Alphabet dates: write down every letter of the alphabet and find a
corresponding date idea for each one.

26 dates, 52 weeks --- one date every two weeks isn't too much to ask.

The restriction adds a bit of creativity, and sometimes a bit of
creative stretching was used to link a date we already wanted to do to a
still-available letter.

This is a goal I can wholeheartedly recommend trying.

\textbf{Success?} YES!\\
\textbf{Was this a good goal?} YES!\\
\textbf{Will it be on the list for 2026?} Probably not --- but something
similar will likely replace it.

\textbf{Climb a mountain}

An adventurous goal that balances out exploring the 10 Belgian
provinces.

Where would I go to find a mountain?

\textbf{Success?} Yes\\
\textbf{Was this a good goal?} Yes\\
\textbf{Will it be on the list for 2026?} Maybe --- I liked the
simplicity and flexibility it offered.

\textbf{Go camping}

This goal also nudges towards simple, cheap, local, adventurous holiday
moments.

And the fact that I didn't achieve this one? I honestly don't have a
good excuse for it.

\textbf{Success?} No\\
\textbf{Was this a good goal?} Yes\\
\textbf{Will it be on the list for 2026?} Yes

\textbf{Moisturize for a month}

This one was only on the list to fill a spot. It was a good resolution
I'd heard from someone else, but one I personally didn't care much
about.

I don't really notice any difference myself.\\
Maybe a month still isn't long enough?

It was a simple goal with low impact --- both positively and negatively.

\textbf{Success?} Yes\\
\textbf{Was this a good goal?} Don't know\\
\textbf{Will it be on the list for 2026?} No

\textbf{Install fly screens}

This was a small renovation we wanted to complete in 2025 anyway.\\
If I have to put in the effort regardless, I might as well let it bring
me closer to my Bingoal.

Because it was on the list, I thought about it more often --- and it may
have gotten done a bit faster because of that.

\textbf{Success?} Yes\\
\textbf{Was this a good goal?} Yes\\
\textbf{Will it be on the list for 2026?} No --- double fly screens
would be weird. But a similar goal might be.

\textbf{Watch 10 ``real'' films}

I'm not very good at sitting still long enough to watch a movie
attentively.\\
I find it easier to read or tell stories than to consume them as films.
But on the other hand, there \emph{are} genuinely great films, and I'm
glad I've seen them --- and extending that list was fun.

\textbf{Success?} Yes\\
\textbf{Was this a good goal?} Yes (for me)\\
\textbf{Will it be on the list for 2026?} Probably. I might raise it to
12 --- one movie per month.

\textbf{Lose 5 kg}

Was this on the list because I wanted to lose weight? Not really.

It was on the list because I'd never actually \emph{tried} to lose
weight and was curious about the experience. During the year, though, I
realized I didn't really care --- so I never truly did anything for it.

I did lose some kilos during the month without cola.\\
Whether they've already come back? No idea --- but after the Christmas
holidays, definitely yes.

\textbf{Success?} No\\
\textbf{Was this a good goal?} No --- not for me\\
\textbf{Will it be on the list for 2026?} No

\textbf{Run or play 20 sessions with the Resonance Protector's Guild
(gaming group)}

This seemed like a realistic --- even slightly pessimistic --- goal for
running roleplaying sessions.

My baseline assumption was one monthly session with a fixed group, plus
one monthly one-shot in the demo group.

The goal was already reached in September, and by the end of the year it
was exceeded by a wide margin.

Still, I'm glad it was on the list. Just like the reading challenge,
this is something I had already been doing for a while without
formalizing it --- and it definitely deserves a place on this list.

\textbf{Success?} Yes\\
\textbf{Was this a good goal?} Yes\\
\textbf{Will it be on the list for 2026?} Yes

\textbf{100 days of Duolingo French}

This was a goal that wasn't on the list yet on January 1st.

But if you live together with someone who has this as a New Year's
resolution, it's an easy one to copy in order to fill another empty
spot.

I actually enjoyed doing this. It helped me speak French a bit more
confidently and fluently during trips to Wallonia. But now, several
months later, that knowledge has completely faded again --- simply
because I don't really use it regularly.

\textbf{Success?} Yes\\
\textbf{Was this a good goal?} Sort of. Languages need to be used and
maintained, so this was largely wasted effort.\\
\textbf{Will it be on the list for 2026?} No --- but another
language-based goal will be.

\textbf{Do a stereotypical group activity (karting, paintball, etc.)}

No challenge is greater than bringing together a group of adults to do
something together.

This challenge was gifted to me for my birthday. My girlfriend had
actually organized a few laser-shooting sessions with some friends ---
which is definitely enough to tick this one off.

\textbf{Success?} Yes\\
\textbf{Was this a good goal?} Yes\\
\textbf{Will it be on the list for 2026?} Maybe

\end{document}
